\section{Kết luận}

Báo cáo nghiên cứu đã phân tích và cải tiến thực nghiệm hệ thống truyền thông tại Trung tâm Công nghệ Giáo dục (daotao.ai). Qua quá trình thực hiện, một số kết luận cốt lõi được rút ra như sau:

\subsection{Tổng kết kết quả thực nghiệm}

Chương trình thử nghiệm thực tế kéo dài hai tuần trên lớp học gồm 100 học viên Đề án 06 đã minh chứng hiệu quả của các giải pháp đề xuất:
\begin{itemize}
    \item Chatbot tương tác trên website đã giải quyết tự động phần lớn yêu cầu kỹ thuật thông thường, trực tiếp giảm số cuộc gọi hỗ trợ trực tiếp từ 24 cuộc xuống còn 4 cuộc, giúp chuyên viên tập trung vào các công việc vận hành chuyên môn khác.
    \item Mẫu email trực quan tinh giản thông tin giúp nâng tỷ lệ tự đăng nhập thành công của học viên trong ngày đầu tiên lên 90\%, tăng đáng kể so với mức 64\% của nhóm đối chứng.
    \item Quy trình Kanban nội bộ giúp tối ưu hóa sự phối hợp kỹ thuật với giảng viên, hạn chế hiện tượng trôi tin nhắn và thiết kế sai học liệu.
\end{itemize}

\subsection{Giá trị của lý thuyết truyền thông trong đào tạo trực tuyến}

Nghiên cứu khẳng định việc ứng dụng các lý thuyết truyền thông kinh điển và hiện đại mang lại giá trị thiết thực để chẩn đoán và thiết kế hệ thống hỗ trợ:
\begin{itemize}
    \item Mô hình truyền thông Shannon–Weaver \parencite{shannon1949} và thuyết phong phú truyền thông \parencite{daft1986} giúp định vị đúng kênh truyền và giảm thiểu nhiễu.
    \item Mô hình giao dịch \parencite{barnlund2008} giúp tái cấu trúc dòng phản hồi tức thì.
    \item Thuyết tải nhận thức \parencite{sweller1988} thiết lập nguyên tắc cốt lõi về sự trực quan và tối giản thông điệp.
\end{itemize}

Sự kết hợp các lý thuyết này tạo ra một khung tư duy hệ thống, giúp thiết kế các điểm chạm truyền thông thân thiện hơn với người học.

\subsection{Đề xuất hướng phát triển tiếp theo}

Để tiếp tục hoàn thiện hệ thống truyền thông của daotao.ai, Trung tâm Công nghệ Giáo dục nên triển khai các bước:
\begin{enumerate}
    \item \textbf{Nâng cấp trí tuệ nhân tạo cho chatbot:} Tiếp tục huấn luyện mô hình ngôn ngữ lớn bằng cơ sở dữ liệu hội thoại mới để nâng tỷ lệ tự động giải quyết và cá nhân hóa lộ trình trả lời.
    \item \textbf{Chuẩn hóa hệ thống đa kênh:} Đồng bộ hóa cơ sở dữ liệu phản hồi giữa chatbot website và các kênh nhắn tin Zalo hay Fanpage để chuyên viên có thể tiếp quản hội thoại đồng bộ mà không làm gián đoạn trải nghiệm của học viên.
    \item \textbf{Đánh giá định kỳ tải nhận thức:} Thực hiện các cuộc đánh giá định kỳ đối với giao diện người dùng và nội dung email thông báo mỗi khi cập nhật tính năng mới trên daotao.ai.
\end{enumerate}
